\chapter{Trabalhos Relacionados}
\label{cap:trabalhos-relacionados}

Este capítulo apresenta trabalhos relacionados ao tema desta pesquisa, com foco em estudos que investigam modularidade, 
qualidade arquitetural e organização de sistemas orientados a objetos. A análise desses trabalhos tem como objetivo identificar
abordagens utilizadas na literatura para avaliação da qualidade estrutural de sistemas, bem como compreender como diferentes
estilos arquiteturais impactam atributos relacionados à manutenção, evolução e organização do software.

Os trabalhos selecionados abordam diferentes perspectivas sobre o problema investigado. Alguns concentram-se na aplicação de métricas 
estruturais para avaliação da qualidade de sistemas orientados a objetos, enquanto outros analisam os impactos de diferentes arquiteturas 
de software em atributos como escalabilidade, manutenibilidade e acoplamento. Entretanto, observa-se que muitos desses estudos são 
fundamentados predominantemente em revisões bibliográficas ou análises conceituais, havendo menor quantidade de pesquisas empíricas 
voltadas à aplicação quantitativa de métricas estruturais em sistemas reais desenvolvidos sob diferentes estilos arquiteturais.

Nesse contexto, os trabalhos apresentados a seguir contribuem para fundamentar teoricamente esta pesquisa e auxiliar na identificação 
de lacunas existentes na literatura, especialmente no que se refere à avaliação quantitativa da modularidade em arquiteturas distintas.

\section{Avaliação de Conectividade em Sistemas Orientados por Objetos}
\label{sec:avaliacao-de-conectividade-em-sistemas-orientados-por-objetos}

O trabalho de \citeonline{ferreira2006} propõe o Modelo de Avaliação de Conectividade de Sistemas Orientados por Objetos (MACSOO), 
fundamentado na tese de que a conectividade entre módulos é o fator preponderante para determinar a estabilidade e a manutenibilidade 
de um software. A pesquisa parte do princípio de que a manutenção consome a maior parte do custo total de um sistema e que estruturas 
mais flexíveis, com menor interdependência entre suas partes, facilitam esse processo.

Segundo a autora, a orientação por objetos, embora favoreça a modularidade, não garante por si só uma estrutura de alta qualidade se não 
houver um controle rigoroso sobre as conexões entre classes. O estudo defende que o impacto de uma alteração no sistema é diretamente 
proporcional ao seu grau de conectividade, e que o "apodrecimento" do design de software ocorre quando a rigidez estrutural impede 
modificações sem gerar efeitos colaterais em cascata.

Para operacionalizar a avaliação, o trabalho identifica e classifica métricas estruturais clássicas, com destaque para o conjunto 
CK (Chidamber e Kemerer) e o conjunto MOOD (Metrics for Object Oriented Design), selecionando indicadores como CBO 
(\textit{Coupling Between Object Classes}), LCOM (\textit{Lack of Cohesion in Methods}) e COF (\textit{Coupling Factor}). 
Além disso, a pesquisa contribuiu com o desenvolvimento da ferramenta \textit{Connecta}, capaz de extrair essas métricas de sistemas 
desenvolvidos em Java e aplicar uma adaptação do modelo de estabilidade de Myers para prever o esforço de manutenção.

A validação do modelo proposto por \citeonline{ferreira2006} foi realizada por meio de experimentos em sistemas reais e acadêmicos, 
demonstrando empiricamente que a dificuldade de manutenção cresce à medida que o acoplamento aumenta e a coesão diminui. 
Entretanto, embora o estudo forneça uma base sólida para a análise de conectividade em sistemas orientados por objetos de forma geral, 
ele não se aprofunda na comparação de como estilos arquiteturais específicos, como o em camadas ou o hexagonal, moldam esses indicadores 
estruturais.

Nesse contexto, o presente trabalho fundamenta-se nas premissas metodológicas de \citeonline{ferreira2006} ao adotar métricas estruturais 
consagradas para avaliar a modularidade de software. Esta pesquisa diferencia-se e avança em relação ao trabalho relacionado ao aplicar 
essa análise quantitativa em um estudo comparativo direto entre a arquitetura em camadas e a arquitetura hexagonal. Enquanto Ferreira 
concentra-se na conectividade como um indicador de qualidade para sistemas OO genéricos, este estudo busca investigar como a escolha 
entre esses dois estilos arquiteturais específicos influencia o acoplamento e a coesão.

\section{Os efeitos da arquitetura dirigida a eventos na modularidade de software}
\label{sec:os-efeitos-da-arquitetura-dirigida-a-eventos-na-modularidade-de-software}

O trabalho de \citeonline{lazzari2022} apresenta um estudo empírico exploratório que busca compreender os efeitos da arquitetura dirigida 
a eventos na modularização de software em comparação ao estilo arquitetural REST. A pesquisa fundamenta-se na análise de duas versões 
funcionalmente idênticas de uma aplicação real, evoluídas através de cinco cenários de mudança nos quais novas funcionalidades foram 
adicionadas. O estudo utiliza um conjunto de dez métricas estruturais para quantificar variáveis críticas da modularidade, 
incluindo separação de interesses (SoC), acoplamento, coesão, complexidade e tamanho.

Segundo os resultados reportados pelos autores, a arquitetura dirigida a eventos (implementada com Apache Kafka) demonstrou 
superioridade na separação de interesses, apresentando melhores índices de modularização funcional tanto em nível de componentes 
quanto de linhas de código. Esse achado sugere que o modelo de eventos favorece a independência dos serviços, facilitando a evolução
do sistema sem a degradação da qualidade funcional. Por outro lado, a pesquisa revelou que essa flexibilidade introduz uma "penalidade" 
estrutural: a arquitetura dirigida a eventos apresentou maior acoplamento e complexidade se comparada ao REST, além de um aumento 
significativo no tamanho do código devido às necessidades de configuração de infraestrutura e classes auxiliares.

O trabalho contribui para a área ao prover evidências quantitativas sobre os trade-offs de um estilo arquitetural moderno em cenários reais 
de evolução. No entanto, os próprios autores reconhecem limitações no estudo, como a análise centrada em uma única aplicação alvo e a elevada 
curva de aprendizado exigida pela arquitetura de eventos, o que pode influenciar a implementação e a extração das métricas. Dessa forma, 
embora o trabalho apresente uma base experimental robusta, os resultados focam especificamente na dicotomia entre comunicação síncrona e 
assíncrona, sem aprofundar-se em estilos organizacionais internos do monólito.

Nesse contexto, o presente trabalho aproxima-se da proposta de \citeonline{lazzari2022} ao adotar uma estratégia de pesquisa experimental 
baseada na comparação de arquiteturas através de métricas estruturais. No entanto, esta pesquisa diferencia-se por 
investigar o impacto de estilos arquiteturais voltados à organização interna do software em vez de focar na comunicação dirigida a eventos.

\section{Comparação entre as arquiteturas de software monolítica e de microsserviços }
\label{sec:comparacao-entre-as-arquiteturas-de-software-monolitica-e-de-microsservicos}

O trabalho de \citeonline{vale2025} investiga o debate entre as arquiteturas monolítica e de microsserviços, focando na transição da 
análise teórica para a validação empírica. O estudo utiliza uma pesquisa quantitativa de natureza experimental para mensurar os trade-offs
de cada abordagem em um ambiente controlado. A metodologia consistiu na implementação de dois protótipos funcionalmente idênticos de um 
sistema de e-commerce (CRUD), desenvolvidos com o ecossistema Node.js e containerizados com Docker, para garantir que as diferenças 
observadas fossem atribuídas exclusivamente às variações arquiteturais.

Segundo os resultados reportados pelo autor, a arquitetura monolítica demonstrou ser objetivamente mais eficiente em todas as métricas de
performance local e custo operacional. O achado mais significativo revelou uma penalidade de desempenho de 147\% na latência de operações 
que exigiam orquestração entre múltiplos domínios na arquitetura de microsserviços, em comparação ao monólito. Além disso, o estudo 
quantificou o "imposto de memória" da distribuição, apontando que os microsserviços consumiram 271\% mais memória RAM em estado de repouso 
devido à sobrecarga de manter múltiplos processos e contêineres ativos.

No que tange à complexidade, o trabalho apresenta uma análise qualitativa indicando que a arquitetura de microsserviços impõe uma 
carga cognitiva e operacional superior. Isso foi evidenciado pelo número de artefatos de configuração necessários para a implantação, 
que foi três vezes maior do que o exigido para o monólito. O autor conclui que, embora os microsserviços ofereçam benefícios teóricos 
de escalabilidade, para sistemas de baixa complexidade ou em fases iniciais, a arquitetura monolítica é fortemente suportada pela eficiência
de recursos e velocidade de entrega.

Embora este TCC compartilhe com \citeonline{vale2025} o uso de uma abordagem quantitativa para analisar trade-offs arquiteturais,
a estratégia de comparação diverge em sua fundamentação. Enquanto o referido autor utiliza protótipos funcionalmente idênticos para 
isolar variáveis de desempenho como latência e consumo de hardware, esta pesquisa foca na análise da modularidade interna por meio de 
métricas estruturais aplicadas a sistemas reais com funcionalidades distintas. Assim, enquanto Vale (2025) quantifica 
custos de infraestrutura em modelos distribuídos, este estudo busca identificar como diferentes padrões organizacionais 
moldam os indicadores de acoplamento e coesão, independentemente da paridade funcional, visando avaliar a sustentabilidade e a 
estabilidade de cada design.

\section{Análise comparativa entre arquiteturas de software: monolítica, modular e microsserviços}
\label{sec:analise-comparativa-entre-arquiteturas-de-software-monolitica-modular-e-microsservicos}
O trabalho de \citeonline{cintra2025} realiza uma análise comparativa entre as arquiteturas monolítica, monólito modular 
e microsserviços, buscando identificar características, vantagens, limitações e cenários de aplicação de cada abordagem arquitetural. 
O estudo é fundamentado em revisão bibliográfica e análise qualitativa, considerando atributos relevantes da Engenharia de Software, 
como escalabilidade, manutenibilidade, resiliência e complexidade operacional.

Segundo os autores, a arquitetura monolítica apresenta vantagens relacionadas à simplicidade de desenvolvimento, implantação e gerenciamento, 
especialmente em sistemas de pequeno e médio porte. Entretanto, à medida que o sistema cresce, problemas relacionados ao aumento do acoplamento
e à dificuldade de manutenção tornam-se mais evidentes. Em contrapartida, a arquitetura de microsserviços oferece maior flexibilidade, 
escalabilidade independente e isolamento de falhas, embora introduza desafios associados à comunicação distribuída, monitoramento e maior 
complexidade operacional \cite{cintra2025}.

O trabalho também destaca o monólito modular como uma alternativa intermediária entre os dois extremos arquiteturais. 
De acordo com os autores, essa abordagem busca manter a simplicidade estrutural do monólito tradicional, porém com maior organização interna 
e separação de responsabilidades, favorecendo a manutenibilidade e facilitando uma possível migração incremental para microsserviços. 
Além disso, o estudo propõe um guia heurístico para apoio à tomada de decisão arquitetural, utilizando pesos e pontuações atribuídos a 
diferentes atributos de qualidade.

Apesar das contribuições apresentadas, os próprios autores reconhecem limitações importantes na pesquisa. O estudo foi conduzido 
predominantemente com base em revisão bibliográfica e análise qualitativa, não sendo realizados experimentos empíricos ou medições 
quantitativas em sistemas reais \cite{cintra2025}. Dessa forma, embora o trabalho apresente uma discussão abrangente sobre 
os impactos das arquiteturas analisadas, não há validação quantitativa baseada em métricas estruturais extraídas diretamente de implementações 
concretas.

Nesse contexto, o presente trabalho aproxima-se da proposta de \citeonline{cintra2025} ao investigar como diferentes estilos 
arquiteturais influenciam atributos relacionados à qualidade estrutural do software. No entanto, esta pesquisa diferencia-se por adotar uma 
abordagem quantitativa baseada na aplicação de métricas estruturais em sistemas reais desenvolvidos sob diferentes arquiteturas. 
Enquanto o trabalho relacionado concentra-se em análises conceituais e heurísticas sobre arquiteturas monolíticas, modulares e distribuídas, 
este estudo busca avaliar empiricamente aspectos relacionados à modularidade, como acoplamento, coesão e estabilidade estrutural, 
comparando arquiteturas em camadas e arquitetura hexagonal.

\section{Um estudo sobre a relação entre qualidade e arquitetura de software}
\label{sec:um-estudo-sobre-a-relacao-entre-qualidade-e-arquitetura-de-software}

O trabalho de \citeonline{tsuruta2010} investiga como a arquitetura de software atua como fator determinante para a qualidade final do 
produto, analisando quatro abordagens distintas para sua especificação: clássica, orientada a objetos, orientada a atributos e orientada 
a busca. A pesquisa fundamenta-se na premissa de que a qualidade não reside apenas na ausência de defeitos funcionais, mas na capacidade do 
sistema de satisfazer requisitos não-funcionais críticos, como a manutenibilidade e a capacidade de modificação.

O autor destaca que o design arquitetural deve ser guiado pelos atributos de qualidade mais importantes para o sucesso do projeto. 
Para viabilizar essa análise, o estudo utiliza a métrica de arquitetura denominada Qualidade de Modularização (TurboMQ), que quantifica
a interdependência entre componentes ao calcular a razão entre relacionamentos internos e externos em um gráfico de dependência de módulos 
(MDG). Um dos diferenciais da pesquisa é a aplicação da Engenharia de Software Baseada em Busca (SBSE), utilizando Algoritmos Genéticos para 
encontrar automaticamente o particionamento ideal de responsabilidades em subsistemas que maximize a modularidade.

Embora forneça uma base teórica robusta sobre modelos de qualidade (como a ISO 9126) e demonstre a eficácia de técnicas meta-heurísticas
na otimização estrutural, o estudo de \citeonline{tsuruta2010} concentra-se em uma aplicação experimental com responsabilidades hipotéticas
para validar o algoritmo. Dessa forma, a pesquisa não realiza uma comparação empírica direta entre padrões arquiteturais consolidados, 
como o estilo em camadas ou o hexagonal, em implementações de software completas. Sob a ótica deste TCC, a investigação aqui conduzida 
compartilha com \citeonline{tsuruta2010} a visão de que a organização dos componentes em subsistemas (ou camadas) é a chave para a 
estabilidade estrutural. Entretanto, este estudo diverge ao adotar uma análise empírica baseada na extração de métricas de sistemas 
reais em vez de buscar a otimização automatizada por algoritmos de busca. Enquanto o referido autor explora o TurboMQ 
para particionamento lógico, esta pesquisa utiliza métricas de acoplamento e coesão para avaliar como os estilos Hexagonal e em Camadas
moldam a modularidade do código-fonte.

\section{Síntese e lacunas}
\label{sec:sintese-e-lacunas}

Após a análise individual dos trabalhos, esta seção apresenta uma visão consolidada das pesquisas selecionadas. 
A Tabela 1 sistematiza os principais atributos e metodologias, facilitando a identificação de lacunas na literatura.

\begin{table}[htb]
\centering
\caption{Comparativo de abordagens dos trabalhos relacionados}
\label{tab:sintese-trabalhos-relacionados}
\begin{tabular}{|l|p{3.5cm}|p{3cm}|p{2.5cm}|c|}
\hline
\textbf{Autor (Ano)} & \textbf{Objetivo Principal} & \textbf{Metodologia} & \textbf{Estilos Arquiteturais} & \textbf{Foco na Mod.} \\ \hline
Ferreira (2006) & Propor o modelo MACSOO para avaliar conectividade. & Experimental (Métricas CK/MOOD em Java). & Sistemas OO genéricos. & Alto \\ \hline
Tsuruta (2010) & Analisar a relação entre qualidade e arquitetura. & Otimização via Busca (Algoritmos Genéticos). & Responsabilidades hipotéticas (TurboMQ). & Médio \\ \hline
Lazzari (2022) & Investigar efeitos da EDA na modularização. & Experimental (Cenários de evolução). & Event-Driven vs. REST. & Alto \\ \hline
Cintra (2025) & Identificar trade-offs para suporte à decisão. & Revisão Bibliográfica Sistemática. & Monolítica, Modular e Microsserviços. & Médio \\ \hline
Vale (2025) & Quantificar custos de desempenho e infraestrutura. & Experimental (Protótipos equivalentes). & Monolítica vs. Microsserviços. & Baixo \\ \hline
\textbf{Este Trabalho} & \textbf{Avaliar modularidade via métricas estruturais.} & \textbf{Experimental (Sistemas SISO e Roomie).} & \textbf{Camadas vs. Hexagonal.} & \textbf{Alto} \\ \hline
\end{tabular}
\vspace{0.1cm}
\par
\centering{\footnotesize Fonte: Elaborado pelo autor (2026).}
\end{table}

A revisão dos trabalhos relacionados demonstra que a análise da qualidade de software tem migrado de discussões puramente teóricas para 
abordagens empíricas fundamentadas em métricas quantitativas. A literatura analisada revela uma forte tendência na utilização de experimentos 
controlados para validar os trade-offs entre diferentes estilos arquiteturais, como observado nos estudos comparativos entre monólitos e 
microsserviços. No que diz respeito à modularidade interna, as métricas estruturais consagradas, como o conjunto CK e MOOD, permanecem como 
os principais instrumentos para quantificar o acoplamento, a coesão e a conectividade de sistemas orientados a objetos.

Entretanto, apesar das contribuições apresentadas pelos trabalhos analisados, ainda existem lacunas importantes relacionadas à
avaliação quantitativa da modularidade em diferentes estilos arquiteturais. A maior parte das investigações comparativas concentra-se 
em arquiteturas de macro-nível, como sistemas monolíticos e distribuídos \cite{cintra2025, vale2025}, ou em diferentes estilos de 
comunicação, como REST e EDA \cite{lazzari2022}. Há uma carência de estudos que comparem de forma quantitativa os padrões
de organização interna de sistemas monolíticos, especificamente a arquitetura em camadas com a arquitetura hexagonal.

Em segundo lugar, identifica-se uma predominância de avaliações qualitativas ou baseadas em revisões bibliográficas que carecem de validação 
em implementações de software concretas \cite{cintra2025}. Mesmo trabalhos que exploram a otimização de atributos de qualidade, como o de 
\citeonline{tsuruta2010}, frequentemente aplicam algoritmos a modelos hipotéticos ou responsabilidades isoladas, sem considerar a complexidade 
de sistemas completos. Além disso, enquanto alguns autores focam em métricas de infraestrutura como latência e consumo de hardware 
\cite{vale2025}, a análise estrutural do código e seu impacto na modularidade é menos explorada em contextos comparativos diretos.

Este TCC propõe-se a preencher essas lacunas ao realizar um estudo experimental focado na modularidade estrutural de dois sistemas: 
o SISO (Arquitetura em Camadas) e o Roomie (Arquitetura Hexagonal). Diferentemente de abordagens que exigem paridade funcional 
entre os sistemas analisados, esta pesquisa aproxima-se da estratégia utilizada por \citeonline{ferreira2006}, avaliando a qualidade 
do design por meio de métricas estruturais aplicadas a sistemas com funcionalidades distintas. Dessa forma, busca-se analisar 
como diferentes estilos arquiteturais influenciam atributos relacionados ao acoplamento e à coesão, produzindo evidências empíricas sobre 
os impactos dessas arquiteturas na modularidade de sistemas orientados a objetos.